\documentclass[12pt]{article}
\usepackage{latexblog}

% ============================================================
% Blog Metadata
% ============================================================
\blogtitle{使用 Antlr 解析配置文件}
\blogdate{2018-05-09}
\blogtags{小实现, 闲话编程, Antlr}
\bloglang{zh}
\blogsummary{}

\begin{document}
\makeblogtitle

在纠结了一阵子 yml,ini,xml甚至 lua 等等 配置文件的格式后，还是决定使用antlr实现了一种我自定义的格式的解析。

这个格式是这个样子的:

\begin{Shaded}
\begin{Highlighting}[]
\OperatorTok{//} \VariableTok{Here} \VariableTok{is} \VariableTok{some} \VariableTok{comment}
\NormalTok{shared }\OperatorTok{\{}
    \VariableTok{string} \VariableTok{\_baseUrl} \OperatorTok{=} \StringTok{"http://localhost:8080"}\OperatorTok{;}
    \VariableTok{string} \VariableTok{domain}   \OperatorTok{=} \StringTok{"riguz.com"}\OperatorTok{;}
    \VariableTok{bool} \VariableTok{ssl}        \OperatorTok{=} \KeywordTok{false}\OperatorTok{;}
    \VariableTok{int} \VariableTok{version}     \OperatorTok{=} \DecValTok{19}\OperatorTok{;}
    \VariableTok{int} \VariableTok{subVersion}  \OperatorTok{=} \DecValTok{25}\OperatorTok{;}
    \VariableTok{float} \VariableTok{number}  \OperatorTok{=} \DecValTok{19.25}\OperatorTok{;}

    \VariableTok{string} \VariableTok{urls}         \OperatorTok{=} \OperatorTok{[}\StringTok{"http://localhost:8080"}\OperatorTok{,} \StringTok{"http://riguz.com:8080"}\OperatorTok{];}
    \VariableTok{string} \VariableTok{domains}      \OperatorTok{=} \OperatorTok{[}\StringTok{"riguz.com"}\OperatorTok{,} \StringTok{"dr.riguz.com"}\OperatorTok{];}
    \VariableTok{bool} \VariableTok{sslArray}       \OperatorTok{=} \OperatorTok{[}\KeywordTok{true}\OperatorTok{,} \KeywordTok{false}\OperatorTok{];}
    \VariableTok{int} \VariableTok{versionArray}    \OperatorTok{=} \OperatorTok{[}\DecValTok{19}\OperatorTok{,} \DecValTok{25}\OperatorTok{];}
    \VariableTok{float} \VariableTok{numberArray}   \OperatorTok{=} \OperatorTok{[}\DecValTok{18.01}\OperatorTok{,} \DecValTok{19.25}\OperatorTok{,} \DecValTok{20.23}\OperatorTok{];}
\OperatorTok{\};}

\VariableTok{scope}\NormalTok{ dev\_db }\OperatorTok{\{}
    \VariableTok{string} \VariableTok{url} \OperatorTok{=}\NormalTok{ $}\OperatorTok{\{}\VariableTok{domain}\OperatorTok{\}} \OperatorTok{..} \StringTok{":3306/mysql"}\OperatorTok{;}
    \VariableTok{string} \VariableTok{user} \OperatorTok{=} \StringTok{"lihaifeng"}\OperatorTok{;}
    \VariableTok{int} \VariableTok{connections} \OperatorTok{=} \DecValTok{10}\OperatorTok{;}
    \VariableTok{string} \VariableTok{password} \OperatorTok{=} \StringTok{"iikjouioqueyjkajkqq=="}\OperatorTok{;}
    \VariableTok{string} \VariableTok{domains} \OperatorTok{=}\NormalTok{ $}\OperatorTok{\{}\VariableTok{domains}\OperatorTok{\};}
\OperatorTok{\};}
\end{Highlighting}
\end{Shaded}

其实是一个k-v形式的文本文件，支持的基本类型有：字符串、布尔值、整数、小数、数组。定义的方法类似于Java或者C语言，

\begin{verbatim}
string _baseUrl = "http://localhost:8080"
\end{verbatim}

前面会限定数据类型。如果要定义数组，则用

\begin{verbatim}
bool sslArray       = [true, false];
\end{verbatim}

这种形式。

然后使用scope区分不同的配置块。因为可能有些相同的配置会重名，这样我们利用不同的scope去区分就好了。考虑到有些配置中需要共同的变量的使用，所以定义了一个shared的scope，这个是写死的scope，其他scope中只能引用shared scope中的变量。

字符串连接使用\texttt{..}操作符。这样可以组装字符串。详细的实现可以在\href{https://github.com/soleverlee/forks/tree/master/config/src/main}{forks的子项目config}中找到。

另外还实现了一个类似Play! Framework的路由定义文件的解析，长这个样子的:

\begin{Shaded}
\begin{Highlighting}[]
\VariableTok{controllers}\NormalTok{ admin}\OperatorTok{\{}
\VariableTok{package} \VariableTok{com}\OperatorTok{.}\VariableTok{riguz}\OperatorTok{.}\VariableTok{forks}\OperatorTok{.}\VariableTok{demo}\OperatorTok{.}\VariableTok{controller}
\VariableTok{UserController}
\VariableTok{FileController}
\OperatorTok{\}}

\NormalTok{controllers }\OperatorTok{\{}
\VariableTok{package} \VariableTok{com}\OperatorTok{.}\VariableTok{riguz}\OperatorTok{.}\VariableTok{forks}\OperatorTok{.}\VariableTok{demo}\OperatorTok{.}\VariableTok{admin}
\VariableTok{UserController}\OperatorTok{{-}\textgreater{}}\VariableTok{AdminUserController}
\VariableTok{PostController}
\OperatorTok{\}}

\NormalTok{filters }\OperatorTok{\{}
\VariableTok{package} \VariableTok{com}\OperatorTok{.}\VariableTok{riguz}\OperatorTok{.}\VariableTok{forks}\OperatorTok{.}\VariableTok{demo}\OperatorTok{.}\VariableTok{filters}
\VariableTok{AuthorizationFilter}
\VariableTok{NocsrfFilter}
\OperatorTok{\}}

\VariableTok{routes}\NormalTok{ admin }\OperatorTok{\{}
\OperatorTok{+}\VariableTok{AuthorizationFilter}
\VariableTok{get}  \OperatorTok{/}\VariableTok{users}                 \VariableTok{UserController}\OperatorTok{.}\NormalTok{getUsers}\OperatorTok{()}
\VariableTok{get}  \OperatorTok{/}\VariableTok{users}\OperatorTok{/:}\VariableTok{id}             \VariableTok{UserController}\OperatorTok{.}\NormalTok{getUser}\OperatorTok{(}\VariableTok{id}\OperatorTok{:} \VariableTok{Long}\OperatorTok{)}
\VariableTok{post} \OperatorTok{/}\VariableTok{users}                 \VariableTok{UserController}\OperatorTok{.}\NormalTok{createUser}\OperatorTok{()}
\VariableTok{get}  \OperatorTok{/}\VariableTok{users}\OperatorTok{/:}\VariableTok{id}\OperatorTok{/}\FunctionTok{files}\OperatorTok{/*}\VariableTok{name} \VariableTok{FileUserController}\OperatorTok{.}\NormalTok{getFile}\OperatorTok{(}\VariableTok{id}\OperatorTok{:} \VariableTok{Long}\OperatorTok{,} \VariableTok{name}\OperatorTok{:} \VariableTok{String}\OperatorTok{)}
\OperatorTok{\}}
\VariableTok{routes}\NormalTok{ guest }\OperatorTok{\{}
\OperatorTok{+}\VariableTok{NocsrfFilter}
\VariableTok{get} \OperatorTok{/}\VariableTok{posts}      \VariableTok{PostUserController}\OperatorTok{.}\NormalTok{getPosts}\OperatorTok{()}
\VariableTok{get} \OperatorTok{/}\VariableTok{posts}\OperatorTok{/:}\VariableTok{id}  \VariableTok{PostUserController}\OperatorTok{.}\NormalTok{getPost}\OperatorTok{(}\VariableTok{id}\OperatorTok{:} \VariableTok{String}\OperatorTok{)}
\OperatorTok{\}}

\VariableTok{routes}\NormalTok{ guest }\OperatorTok{\{}
\OperatorTok{+}\VariableTok{NocsrfFilter}
\VariableTok{get} \OperatorTok{/}\VariableTok{posts}      \VariableTok{PostUserController}\OperatorTok{.}\NormalTok{getPosts}\OperatorTok{()}
\VariableTok{get} \OperatorTok{/}\VariableTok{posts}\OperatorTok{/:}\VariableTok{id}  \VariableTok{PostUserController}\OperatorTok{.}\NormalTok{getPost}\OperatorTok{(}\VariableTok{id}\OperatorTok{:} \VariableTok{String}\OperatorTok{)}
\OperatorTok{\}}
\end{Highlighting}
\end{Shaded}

这个文件的解析也在上面的git中可以找到实现。通过Antlr可以很方便的把类似这样的文件解析出来，你甚至可以实现自己的领域语言。在实现过程中，遇到过一些问题，来说下问题吧。

首先是Antlr提供了Listener和Visitor两种方式，起初使用Listener来实现但是感觉比较麻烦，而使用Visitor则可以直接通过返回值来取得AST解析结果。我们解析一个文件的时候，是自顶向下的，一个个的去解析的，比如我们的配置文件的antlr语法定义如下：

\begin{Shaded}
\begin{Highlighting}[]
\VariableTok{script}
    \OperatorTok{:} \VariableTok{shared}\NormalTok{? }\VariableTok{scope}\OperatorTok{*}
      \ConstantTok{EOF}
    \OperatorTok{;}
\end{Highlighting}
\end{Shaded}

其中shared又是这样的

\begin{Shaded}
\begin{Highlighting}[]
\VariableTok{shared}
    \OperatorTok{:} \ConstantTok{SHARED}\NormalTok{ LBRACE }\OperatorTok{(}\VariableTok{property} \ConstantTok{SEMI}\OperatorTok{)*} \ConstantTok{RBRACE} \ConstantTok{SEMI}
    \OperatorTok{;}
\end{Highlighting}
\end{Shaded}

也就是说 \texttt{shared\ \{\ k=v...\}\ ;}这样的形式，然后又开始到了property:

\begin{Shaded}
\begin{Highlighting}[]
\VariableTok{property}
    \OperatorTok{:} \FunctionTok{type} \ConstantTok{NAME} \ConstantTok{ASSIGN} \VariableTok{expression}        \OperatorTok{\#}\VariableTok{basicProperty}
    \OperatorTok{|} \FunctionTok{type} \ConstantTok{NAME} \ConstantTok{ASSIGN} \ConstantTok{LBRACK}
        \VariableTok{expression}\NormalTok{? }\OperatorTok{(}\ConstantTok{COMMA} \VariableTok{expression}\OperatorTok{)*}
      \ConstantTok{RBRACK}                             \OperatorTok{\#}\VariableTok{arrayProperty}
    \OperatorTok{;}
\end{Highlighting}
\end{Shaded}

这样层层往下来看。然后解析的时候也是一样，我们首先有一个顶层的解析器：

\begin{Shaded}
\begin{Highlighting}[]
\KeywordTok{public} \KeywordTok{class}\NormalTok{ ScriptVisitor }\KeywordTok{extends}\NormalTok{ CfParserBaseVisitor}\OperatorTok{\textless{}}\BuiltInTok{Map}\OperatorTok{\textless{}}\BuiltInTok{String}\OperatorTok{,}\NormalTok{ ScriptVisitor}\OperatorTok{.}\FunctionTok{Scope}\OperatorTok{\textgreater{}\textgreater{}} \OperatorTok{\{}
    \KeywordTok{private} \DataTypeTok{static} \DataTypeTok{final} \BuiltInTok{Logger}\NormalTok{ logger }\OperatorTok{=}\NormalTok{ LoggerFactory}\OperatorTok{.}\FunctionTok{getLogger}\OperatorTok{(}\NormalTok{ScriptVisitor}\OperatorTok{.}\FunctionTok{class}\OperatorTok{);}

    \AttributeTok{@Override}
    \KeywordTok{public} \BuiltInTok{Map}\OperatorTok{\textless{}}\BuiltInTok{String}\OperatorTok{,}\NormalTok{ Scope}\OperatorTok{\textgreater{}} \FunctionTok{visitScript}\OperatorTok{(}\NormalTok{CfParser}\OperatorTok{.}\FunctionTok{ScriptContext}\NormalTok{ ctx}\OperatorTok{)} \OperatorTok{\{}
        \KeywordTok{...}
    \OperatorTok{\}}
\end{Highlighting}
\end{Shaded}

这个Visitor负责解析语法文件中定义的script块，然后解析里面的scope：

\begin{Shaded}
\begin{Highlighting}[]
\NormalTok{ScopeVisitor scopeVisitor }\OperatorTok{=} \KeywordTok{new} \FunctionTok{ScopeVisitor}\OperatorTok{(}\NormalTok{context}\OperatorTok{);}
\NormalTok{        ctx}\OperatorTok{.}\FunctionTok{scope}\OperatorTok{().}\FunctionTok{forEach}\OperatorTok{(}\NormalTok{scopeContext }\OperatorTok{{-}\textgreater{}} \OperatorTok{\{}
\NormalTok{            logger}\OperatorTok{.}\FunctionTok{debug}\OperatorTok{(}\StringTok{"Visit scope:\{\}"}\OperatorTok{,}\NormalTok{ scopeContext}\OperatorTok{.}\FunctionTok{getText}\OperatorTok{());}
\NormalTok{            Scope scope }\OperatorTok{=}\NormalTok{ scopeContext}\OperatorTok{.}\FunctionTok{accept}\OperatorTok{(}\NormalTok{scopeVisitor}\OperatorTok{);}
\NormalTok{            scopes}\OperatorTok{.}\FunctionTok{put}\OperatorTok{(}\NormalTok{scope}\OperatorTok{.}\FunctionTok{name}\OperatorTok{,}\NormalTok{ scope}\OperatorTok{);}
        \OperatorTok{\});}
\end{Highlighting}
\end{Shaded}

这样又实现一个ScopeVisitor去解析scope就好了。详细的实现就不多贴代码了。

另外一个问题是，对于错误的处理，我们在哪一步做？比如\texttt{bool\ s\ =\ "123";}这是错误的，我们其实可以在定义grammar的时候就避免这种错误来，但写起来会麻烦一些。目前的实现是在Visitor中去对逻辑进行判断的，前面只做语法检查就可以了。

参考:

\begin{itemize}
\tightlist
\item
  http://jakubdziworski.github.io/java/2016/04/01/antlr\_visitor\_vs\_listener.html
\end{itemize}


\end{document}
